% Page de titre, résumés et mots-clés (français et anglais) et classification JEL de budget.tex
\documentclass[12pt]{article}
\usepackage[T1]{fontenc}
\usepackage[utf8]{inputenc}
\usepackage{tgpagella}
\usepackage{mathpazo}
\usepackage{microtype}
\usepackage[margin=2.5cm]{geometry}
\usepackage{setspace}
\onehalfspacing
\usepackage[english,french]{babel}
\usepackage{csquotes}
\usepackage[colorlinks=true,urlcolor=blue]{hyperref}
\pagestyle{empty}

\begin{document}

\begin{center}
  {\Large\bfseries Préférences budgétaires des Français~:\\ résultats d'une enquête représentative\par}
  \medskip
  {\large Adrien Fabre\par}
\end{center}

\section*{Résumé}
J'analyse les résultats d'une enquête sur les préférences programmatiques et budgétaires menée en 2026 sur un échantillon représentatif de 1\,000 Français. Les Français préfèrent les mesures de rationalisation de la dépense publique, les taxes sur les plus riches et la suppression des aides aux étrangers~; ils rejettent la hausse de l'âge de départ à la retraite et les hausses de taxes indifférenciées (TVA, CSG). Par ailleurs, une majorité souhaite améliorer les services publics, augmenter le SMIC, ou une plus grande fermeté judiciaire. La tripartition des partis entre gauche, centre-droit/droite et extrême-droite émerge naturellement des préférences programmatiques. La gauche se distingue principalement sur la question des étrangers, le centre-droit/droite sur celle des retraites. Le paquet de mesures majoritaire permettant la plus grande réduction du déficit réalise 70~Mds d'économie, à l'aide de taxes accrues sur les plus aisés, du gel des dépenses publiques, et de taxes écologiques sectorielles. Un groupe de \textit{frugaux} qui soutient tout type de mesures de consolidation budgétaire pourrait déterminer la prochaine élection en faisant gagner le candidat au programme budgétaire le plus crédible. Enfin, j'observe que les préférences individuelles sont bien plus variées au sein des électeurs d'un même parti qu'entre préférences moyennes des électeurs de partis différents.

\medskip\noindent
\textbf{Mots-clés~:} Préférences~; déficit~; sondage~; opinion~; attitudes.

\clearpage
\begin{otherlanguage}{english}
\begin{center}
  {\Large\bfseries French Fiscal Preferences:\\ Results from a Representative Survey\par}
\end{center}

\section*{Abstract}
I analyze the results of a survey on programmatic and budgetary preferences conducted in 2026 on a representative sample of 1,000 French people. The French prefer measures that rationalize public spending, taxes on the wealthiest, and the elimination of welfare benefits for foreigners; they reject raising the retirement age and across-the-board tax hikes (VAT, CSG). Moreover, a majority wishes to improve public services, raise the minimum wage, or see tougher law enforcement. The tripartite division of parties between the left, the centre-right/right and the far right emerges naturally from programmatic preferences. The left stands out mainly on the question of foreigners, the centre-right/right on pensions. The majority-backed package of measures achieving the largest deficit reduction yields 70 billion euros of savings, through higher taxes on the wealthiest, a freeze on public spending, and sector-specific environmental taxes. A group of deficit hawks who support any kind of fiscal consolidation measure could determine the next election by ensuring the victory of the candidate with the most credible budget plan. Finally, I observe that individual preferences are far more varied among the voters of a given party than between the average preferences of voters of different parties.

\medskip\noindent
\textbf{Keywords:} Preferences; deficit; survey; opinion; attitudes.
\end{otherlanguage}

\bigskip\noindent
\textbf{Classification JEL~:} D72~; H60~; H63~; P16.

\end{document}
